\documentclass[10pt]{article}

\usepackage[margin=0.75in]{geometry}
\usepackage{amsmath,amsthm,amssymb,bm}
\usepackage{xcolor}
\usepackage{cancel}
\usepackage{graphicx}
\usepackage{changepage}
\usepackage{circuitikz}
\usepackage{pgfplots}
\usepackage{minted}
\usepackage{physics}
\usepackage{tikz}
\usepackage{tikz-3dplot}
\usepackage{hyperref}
\usepackage{cleveref}
\usepackage{siunitx}
\usepackage{fontspec}
\usepackage{relsize}
\usepackage{subfig}
\usepackage{todonotes}
\usepackage{contour}
\usepackage{multicol, multirow, booktabs}
\usepackage[breakable]{tcolorbox}
\usepackage[inline]{enumitem}

\theoremstyle{definition}
\newtheorem{problem}{Problem}
\newtheorem{soln}{Solution}

\pgfplotsset{compat=newest}
\usetikzlibrary{lindenmayersystems}
\usetikzlibrary{arrows}
\usetikzlibrary{calc}
\usetikzlibrary{positioning, fit}
\usetikzlibrary{3d, perspective}
\usetikzlibrary{math}
\usetikzlibrary{fadings}
\usetikzlibrary{angles,quotes} 
\usetikzlibrary{decorations.markings}

\definecolor{incolor}{HTML}{303F9F}
\definecolor{outcolor}{HTML}{D84315}
\definecolor{cellborder}{HTML}{CFCFCF}
\definecolor{cellbackground}{HTML}{F7F7F7}
\newcommand{\ui}{\hat{i}}
\newcommand{\uj}{\hat{j}}
\newcommand{\uk}{\hat{k}}
\newcommand{\ux}{\hat{\mathbf{x}}}
\newcommand{\uy}{\hat{\mathbf{y}}}
\newcommand{\uz}{\hat{\mathbf{z}}}
\newcommand{\uv}[1]{\hat{\bv{{#1}}}}
\newcommand{\pr}[1]{{#1^\prime}}
\newcommand{\pd}[2][]{{\frac{\partial^{#1}}{\partial {{#2}^{#1}}}}}
\newcommand{\dif}[2][]{{\frac{d^{#1}}{d{{#2}^{#1}}}}}
\newcommand{\grd}{\nabla}

\pgfdeclarelayer{background}  
\pgfsetlayers{background,main}
\AtBeginDocument{\RenewCommandCopy\qty\SI}

\makeatletter
\newcommand{\boxspacing}{\kern\kvtcb@left@rule\kern\kvtcb@boxsep}
\makeatother
\newcommand{\prompt}[4]{
    \ttfamily\llap{{\color{#2}[#3]:\hspace{3pt}#4}}\vspace{-\baselineskip}
}

\newcommand{\thevenin}[2]{
  \begin{center}
    \begin{circuitikz} \draw
      (0,0) -- (2,0) to[battery1, l_=$V_{Th}\eq#1$] (2,2) 
      to[resistor, l_=$R_{Th}\eq#2$] (0,2)
      ;
      \draw [o-] (-.07,2.079);
      \draw [o-] (-.07,0.079);
    \end{circuitikz}
  \end{center}
}

\newcommand{\norton}[2]{
  \begin{center}
    \begin{circuitikz} \draw
      (0,0) -- (3,0) to[american current source, l_=$I_{N}\eq#1$] (3,2) -- (0,2) (2,0)
      to[resistor, l=$R_{N}\eq#2$] (2,2)
      ;
      \draw [o-] (-.07,2.079);
      \draw [o-] (-.07,0.079);
    \end{circuitikz}
  \end{center}
}

\DeclareMathOperator{\Div}{div}
\DeclareMathOperator{\Curl}{curl}
\DeclareMathOperator{\sgn}{sgn}

\newcommand{\highlight}[1]{\colorbox{yellow}{$\displaystyle #1$}}
\newcommand{\ti}[1]{\widetilde{#1}}
\renewcommand{\div}[1]{\bv{\nabla}\cdot #1}
\renewcommand{\curl}[1]{\bv{\nabla}\times #1}

\newfontface{\Kaufmann}{Kaufmann}
\DeclareTextFontCommand{\kf}{\Kaufmann}
\newcommand{\scriptr}{\fontsize{12pt}{12pt}\kf{r}\fontsize{10pt}{10pt}}

\newfontface{\KaufmannB}{Kaufmann Bold}
\DeclareTextFontCommand{\kfb}{\KaufmannB}
\newcommand{\bscriptr}{\fontsize{12pt}{12pt}\kfb{r}\fontsize{10pt}{10pt}}
\newcommand{\justif}[2]{&{#1}&\text{#2}}
\newcommand{\bv}[1]{\bm{\mathrm{#1}}}

\title{Physics 4220H: Assignment I}
\author{Jeremy Favro (0805980) \\ Trent University, Peterborough, ON, Canada}
\date{\today}

\begin{document}
\maketitle

% PROBLEM 1
\begin{problem}
Derive the four typical boundary conditions given below from (the complete)
Maxwell's equations. Please include diagrams as needed (drawn by hand is fine):
\begin{align*}
  \bv{E}_1^\parallel-\bv{E}_2^\parallel                               & =0                    \\
  \frac{1}{\mu_1}\bv{B}_1^\parallel-\frac{1}{\mu_2}\bv{B}_2^\parallel & =\bv{K}_f\times\uv{n} \\
  \epsilon_1E_1^\perp-\epsilon_2E_2^\perp                             & =\sigma_f             \\
  B_1^\perp-B_2^\perp                                                 & =0
\end{align*}
\end{problem}
\begin{soln}

\end{soln}

% PROBLEM 2
\begin{problem}
Derive the charge continuity equation,
$$\pd{t}\rho(\bv{r},t)=-\div{\bv{J}(\bv{r},t)}$$
from (the complete) Maxwell's Equations.
\end{problem}
\begin{soln}
  Amp\`ere's law states that
  $$\curl{\bv{B}}=\mu_0\bv{J}+\mu_0\epsilon_0\frac{\partial \bv{E}}{\partial t}.$$
  Taking the divergence of both sides of this we get that
  $$\div{(\curl{\bv{B}})}=0=\mu_0\div{\bv{J}}+\mu_0\epsilon_0\frac{\partial}{\partial t}\div{\bv{E}}$$
  which can be rearranged to
  \begin{equation}
    \div{\bv{J}}=-\epsilon_0\frac{\partial}{\partial t}\div{\bv{E}}.
    \label{eqn:nc}
  \end{equation}
  Gauss' law then states that
  \begin{equation}
    \div{E}=\frac{\rho}{\epsilon_0}.
    \label{eqn:gauss}
  \end{equation}
  Plugging \cref{eqn:gauss} into \cref{eqn:nc} we get the continuity equation,
  $$\div{\bv{J}}=-\frac{\partial\rho}{\partial t}.$$
\end{soln}

% PROBLEM 3
\begin{problem}~
\begin{enumerate}[label=(\alph*)]
  \item Show that, in linear media, the Poynting vector can be given by $\bv{S}=\bv{E}\times\bv{H}$, and that the energy density is
        $u_{EM}=\frac{1}{2}\left(\bv{E}\cdot\bv{D}+\bv{B}\cdot \bv{H}\right)$.
  \item Describe in words the conceptual difference between the Poynting vector and
        the total electromagnetic energy density. Can you give an example of when
        computing one might be more relevant than the other?
\end{enumerate}
\end{problem}
\begin{soln}~
  \begin{enumerate}[label=(\alph*)]
    \item I'll follow basically the same process we did in class except using Maxwell's Equations in linear media,
          \begin{align*}
            \div{\bv{D}}=\rho_f                               \\
            \curl{\bv{E}}=-\frac{\partial \bv{B}}{\partial t} \\
            \div{\bv{B}}=0                                    \\
            \curl{\bv{H}}=\bv{J}_f+\frac{\partial \bv{D}}{\partial t}.
          \end{align*}
          Now the work done by an electromagnetic field on a charge in infinitesimal timestep $dt$ (over which a length $d\ell$ is traveled)
          is $dW=\bv{F}\cdot d\bv{\ell}=q(\bv{E}+\bv{v}\times\bv{B})\cdot d\ell$. Noting that $\bv{v}=\frac{d\bv{\ell}}{dt}$ we can rewrite this as
          $$dW=q(\bv{E}+\bv{v}\times\bv{B})\cdot \bv{v}dt.$$ Noting then that $\bv{v}\times\bv{B}$ by definition of the cross product
          gives a vector perpendicular to $\bv{v}$ and $\bv{B}$, dotting the output with $\bv{v}$ will give zero, hence, dividing through by $dt$,
          $$\frac{dW}{dt}=\bv{E}\cdot q\bv{v}.$$
          Then since $q=\int_\mathcal{V}\rho(\bv{r}')dr'$,
          $$\frac{dW}{dt}=\bv{E}\cdot \int_\mathcal{V}\rho(\bv{r}')\bv{v}dr'=\int_\mathcal{V}\bv{E}\cdot \bv{J}dr'$$
          provided $\bv{E}$ is position independent in $\mathcal{V}$ and since $\rho \bv{v}$ is a charge moving, i.e. a current.
          Now applying Amp\`ere's law in media,
          $$\bv{J}=\curl{\bv{H}}-\frac{\partial \bv{D}}{\partial t},$$
          $$\frac{dW}{dt}=\int_\mathcal{V} \bv{E}\cdot (\curl{\bv{H}})-\bv{E}\cdot \frac{\partial \bv{D}}{\partial t}dr'.$$
          Now to fiddle with the first term we use the vector identity $\div{(\bv{E}\times\bv{H})}=(\curl{\bv{E}})\cdot\bv{H}-\bv{E}\cdot(\curl{\bv{H}})$ and obtain
          $$\frac{dW}{dt}=\int_\mathcal{V} (\curl{\bv{E}})\cdot\bv{H}-\div{(\bv{E}\times\bv{H})}-\bv{E}\cdot \frac{\partial \bv{D}}{\partial t}dr'.$$
          Now applying Faraday's law to the first term,
          $$\frac{dW}{dt}=-\int_\mathcal{V} \bv{H}\cdot\frac{\partial \bv{B}}{\partial t}+\bv{E}\cdot \frac{\partial \bv{D}}{\partial t} + \div{(\bv{E}\times\bv{H})}dr'.$$
          Now using the bit of a priori knowledge from class regarding the chain-rule differentiation trick we can look more closely at the first two terms.
          Since $\bv{B}=\mu_0(\bv{H}+\bv{M})=\mu\bv{H}$ where $\mu=\mu_0(1+\chi_m)$ and 
          $$\frac{1}{\mu}\frac{\partial \abs{\bv{B}}^2}{\partial t}=\frac{2}{\mu}\bv{B}\cdot\frac{\partial \bv{B}}{\partial t}=2\bv{H}\cdot\frac{\partial \bv{B}}{\partial t}$$
          we can rewrite the first term as 
          $$\bv{H}\cdot\frac{\partial \bv{B}}{\partial t}=\frac{1}{2}\frac{\partial \bv{B}\cdot\bv{H}}{\partial t}$$
          and the second term using the same chain rule approach with $\bv{D}=\epsilon\bv{E}=\epsilon_0(1+\chi_e)\bv{E}$ as
          $$\bv{E}\cdot \frac{\partial \bv{D}}{\partial t}=\frac{1}{2} \frac{\partial \bv{E}\cdot\bv{D}}{\partial t}.$$
          The main expression then becomes 
          $$\frac{dW}{dt}=-\int_\mathcal{V} \frac{1}{2}\frac{\partial }{\partial t}\left[\bv{B}\cdot\bv{H}+\bv{E}\cdot\bv{D}\right]dr' - \int_\mathcal{V}\div{(\bv{E}\times\bv{H})}dr'.$$
          We can fiddle with the last term with the divergence theorem (and pull out the time derivative in the first) to obtain
          $$\frac{dW}{dt}=-\frac{1}{2}\frac{d }{d t}\int_\mathcal{V} \left[\bv{B}\cdot\bv{H}+\bv{E}\cdot\bv{D}\right]dr' - \oint_\mathcal{S}(\bv{E}\times\bv{H})\cdot d\bv{a}.$$
          I think at this point we can just look at Poynting's theorem in vacuum and compare the two. Performing such a comparison makes it immediately obvious that the 
          term in the first integral is an analogue to $u_{EM}$, the electromagnetic energy density and the term in the second integral is the Poynting vector.
    \item To my understanding, due to the way the $dW/dt$ equation is written, the $u_{EM}$ term represents the total energy stored in the fields in a volume $\mathcal{V}$. The
    Poynting vector integral term represents the flow of energy through $\mathcal{S}$, the surface which bounds $\mathcal{V}$. \todo{Situation where one is more useful than the other?}
  \end{enumerate}
\end{soln}

% PROBLEM 4
\begin{problem}
The Dirac Quantization Condition is a famous rule in quantum mechanics that states
that if even a single magnetic monopole exists in the universe, electric charge must
be quantized. There is a famously cute (if technically improper) way to “show” this,
involving quantization of field angular momentum:


Suppose that you had an electric point charge, $q_e$, separated by distance $R$ from a
magnetic point charge (or monopole), $q_m$. Put the monopole at the origin, and the
electric charge at $\bv{R}$. The magnetic and electric fields emanating from these charges
would be
$$\bv{B}=\frac{\mu_0}{4\pi}\frac{q_m}{r^3}\bv{r}$$
and
$$\bv{E}=\frac{1}{4\pi\epsilon_0}\frac{q_e}{\abs{\bv{r}-\bv{R}}^3}(\bv{r}-\bv{R}).$$
\begin{enumerate}[label=(\alph*)]
  \item Calculate the field momentum density, $\bv{p}_{EM}$, of this arrangement (known as a
        Thompson Dipole, btw), as a function of $r$, and comment on the direction of this
        momentum density.
  \item From this, find the ``field angular momentum'' possessed by this arrangement, and
        comment on
        \begin{enumerate}[label=\roman*.]
          \item the meaning (or strangeness) of the functional dependence of this angular
                momentum on the distance between the two “charges”,
          \item the direction of this angular momentum,
          \item whether it is intuitive to you that static “charges” can set up fields arranged with
                nonzero angular momentum.
        \end{enumerate}
        [Note, the integral involved is notoriously difficult, usually requiring several
        shortcuts or tricks. If needed, you may invoke the following result:
        $$\int\frac{R}{r}\frac{1}{\abs{\bv{r}-\bv{R}}}\left[\uv{R}-\uv{r}\left(\uv{r}\cdot\uv{R}\right)\right]dV=4\pi \uv{R}]$$
  \item A fundamental requirement of quantum mechanics is that angular momentum be
        quantized in units of $\hbar/2$. Use this fact, together with your result from (b) to argue for
        Dirac quantization.
\end{enumerate}
\end{problem}
\begin{soln}

\end{soln}
\end{document}